\documentclass[11pt]{article}
\usepackage[T1]{fontenc}
\usepackage[utf8]{inputenc}
\usepackage[spanish,es-noshorthands]{babel}
\usepackage{lmodern}
\usepackage[margin=2.4cm]{geometry}
\usepackage{amsmath,amssymb}
\usepackage{booktabs}
\usepackage{xcolor}
\usepackage{enumitem}
\usepackage{parskip}
\usepackage[colorlinks=true,linkcolor=black,urlcolor=blue]{hyperref}

\definecolor{acmuerta}{HTML}{9AA0A6}
\definecolor{aclatente}{HTML}{C9A227}
\definecolor{acactiva}{HTML}{2E7D32}
\newcommand{\MUERTA}{\textcolor{acmuerta}{\textsf{MUERTA}}}
\newcommand{\LATENTE}{\textcolor{aclatente}{\textsf{LATENTE}}}
\newcommand{\ACTIVA}{\textcolor{acactiva}{\textsf{ACTIVA}}}

\title{\textbf{AstroBioSim — Modelo biológico de las especies}\\[2pt]
\large Decisiones de diseño para la exposición y defensa final}
\author{Área de Biotecnología}
\date{Julio 2026}

\begin{document}
\maketitle

\begin{abstract}
\noindent El motor biológico define \emph{quién} vive y \emph{dónde}: tres
microorganismos análogos, cada uno con umbrales de tolerancia calibrados con
literatura. Este documento justifica la elección de especies, la separación entre
\emph{crecimiento} y \emph{supervivencia} (que da lugar a la latencia), la
asignación especie–entorno y la trazabilidad de cada parámetro. La regla de oro:
un parámetro sin procedencia no entra al modelo.
\end{abstract}

\section{Las tres especies análogas}

\begin{center}\small
\begin{tabular}{@{}llll@{}}
\toprule
Entorno & Especie & Rasgo clave & Metabolismo \\
\midrule
Tierra (control) & \emph{Escherichia coli} & mesófila estándar & heterótrofo \\
Marte & \emph{Deinococcus radiodurans} & radiorresistente, anhidrobiótica & heterótrofo \\
Encelado & \emph{Methanococcoides burtonii} & psicrotolerante & metanógena \\
\bottomrule
\end{tabular}
\end{center}

\emph{E.\ coli} es el control: no un extremófilo, sino la referencia con
parámetros de laboratorio bien establecidos. \emph{D.\ radiodurans} es el candidato
marciano por su resistencia a radiación y desecación. \emph{M.\ burtonii}
reemplazó a una termófila original (ADR-0011) porque los datos de ventila dan
$T\approx 2{,}4\,^\circ$C: una termófila moriría trivialmente. \emph{M.\ burtonii}
es psicrotolerante (lago Ace, Antártida) y \textbf{metanógena}, lo que conserva la
relevancia astrobiológica directa —el penacho de Encelado contiene CH$_4$ y H$_2$.

\section{Crecimiento vs.\ supervivencia: el origen de la latencia (ADR-0012)}

La decisión biológica de fondo. Cada especie declara \emph{dos} juegos de umbrales,
no uno:
\begin{itemize}[nosep]
  \item \textbf{Crecimiento} ($\mu>0$): $t_{\min},t_{\mathrm{opt}},t_{\max}$,
  $a_{w,\min}$, $\mathrm{UV}_{\max}$.
  \item \textbf{Supervivencia} (sigue viva, $\mu=0$): $t_{\text{sup,min}},
  t_{\text{sup,max}}$, $a_{w,\text{sup}}$, $\mathrm{UV}_{\text{letal}}$.
\end{itemize}
El motivo es la \emph{anhidrobiosis}: \emph{D.\ radiodurans}, al quedarse sin agua,
no muere sino que se \emph{apaga} y revive con el agua. Su límite de metabolismo
\emph{activo} es $a_w\approx 0{,}90$, pero sobrevive seca hasta $a_w\approx 0$. Un
modelo de un solo umbral obligaría a elegir entre matarla (falso) o declarar
crecimiento por debajo del límite de la vida (indefendible). Con dos umbrales, en
Atacama queda \LATENTE: viva, sin crecer —que es la verdad.

\paragraph{Cota dura.} Ningún $a_{w,\min}$ de crecimiento puede bajar de
$\mathbf{0{,}605}$, el límite inferior de división celular conocido para cualquier
organismo terrestre (Stevenson \emph{et al.}, 2015). Un test lo verifica y falla si
alguien lo cruza. Los tres estados \MUERTA/\LATENTE/\ACTIVA{} del autómata son la
traducción directa de esta biología.

\section{Parámetros calibrados y su procedencia}

Toda cifra lleva etiqueta de procedencia: \textbf{[LIT]} publicado, \textbf{[DER]}
derivado por fórmula, \textbf{[ANA]} inferido por analogía con un organismo medido,
\textbf{[CONV]} convención, \textbf{[EST]} estimación pendiente.

\begin{center}\small
\begin{tabular}{@{}lcccl@{}}
\toprule
Especie & $t_{\min}/t_{\mathrm{opt}}/t_{\max}$ (\textdegree C) & $a_{w,\min}$ & $a_{w,\text{sup}}$ & Proc.\ (agua) \\
\midrule
\emph{E.\ coli} & $7{,}5\,/\,37\,/\,47$ & $0{,}95$ & $0{,}50$ & [LIT] / [EST] \\
\emph{D.\ radiodurans} & $4\,/\,30\,/\,39$ & $0{,}90$ & $0{,}0$ & [LIT] / [LIT] \\
\emph{M.\ burtonii} & $-2{,}5\,/\,23{,}4\,/\,29{,}5$ & $0{,}95$ & $0{,}80$ & [LIT] / [EST] \\
\bottomrule
\end{tabular}
\end{center}

\paragraph{Umbrales de UV, derivados de fluencias (ADR-0014).} La letalidad UV es
una \emph{fluencia} (J/m$^2$) y el campo es una \emph{irradiancia} (W/m$^2$); se
convierten dividiendo por la exposición por tick:
\begin{equation}
\mathrm{UV}_{\text{letal}}=\frac{F_{\text{letal}}}{\tau},\qquad
\mathrm{UV}_{\max}=\frac{\mathrm{UV}_{\text{letal}}}{10},\qquad \tau\approx 8\ \text{h de sol.}
\end{equation}
Con $F_{\text{letal}}=870$~J/m$^2$ para \emph{E.\ coli} y $50\,760$~J/m$^2$ para
\emph{D.\ radiodurans} [LIT]: un factor $\mathbf{58\times}$ que sale del dato, no de
una estimación. Para \emph{M.\ burtonii} no hay medición publicada; se infiere
\textbf{[ANA]}: las metanógenas adaptadas al frío resisten $2{,}5$--$13{,}8\times$
más UV que la referencia sensible (\emph{M.\ barkeri}, $\approx$\emph{E.\ coli}), y
se toma el extremo inferior ($2{,}5\times$) por prudencia.

\section{Habitabilidad como intersección booleana}

Para un campo dado, la máscara de crecimiento de cada especie es la intersección
vectorizada de sus umbrales:
\[
\text{crece}=\big(T\ge t_{\min}\big)\wedge\big(T\le t_{\max}\big)
\wedge\big(a_w\ge a_{w,\min}\big)\wedge\big(\mathrm{UV}\le \mathrm{UV}_{\max}\big),
\]
y análogamente la de supervivencia con los umbrales laxos. La de supervivencia es
siempre un \emph{superconjunto} de la de crecimiento (se testea).

\section{Resultado del modelo y lectura biológica}

Con los parámetros calibrados, cada especie en su entorno da:
\begin{center}\small
\begin{tabular}{@{}lccc@{}}
\toprule
Especie / entorno & \ACTIVA & \LATENTE & \MUERTA \\
\midrule
\emph{E.\ coli} / Tierra & $100\%$ & $0\%$ & $0\%$ \\
\emph{D.\ radiodurans} / Marte & $0\%$ & $98\%$ & $2\%$ \\
\emph{M.\ burtonii} / Encelado & $100\%$ & $0\%$ & $0\%$ \\
\bottomrule
\end{tabular}
\end{center}

La lectura es coherente: el control prospera; \emph{M.\ burtonii} crece incluso en
los núcleos de fumarola (con $t_{\max}=29{,}5$, la vida se agrupa en las ventilas,
no las evita); y \emph{D.\ radiodurans} en Marte \emph{sobrevive latente pero no
crece} —el regolito en bulk es inhabitable, y solo un microrefugio húmedo
reactiva la división. Que Marte dé extinción del crecimiento activo \textbf{no es
un error}: es el resultado correcto y el punto de partida de la pregunta de
investigación.

\section*{Supervisión biológica (criterios de la defensa)}
\begin{itemize}[nosep]
  \item Los umbrales citan literatura; el único \textbf{[EST]} que resta es
  $a_{w,\text{sup}}$ de \emph{E.\ coli} y \emph{M.\ burtonii}.
  \item Se respeta que \emph{ninguna célula muerta revive}: \MUERTA{} es absorbente;
  \LATENTE$\to$\ACTIVA{} sí es válido (no es resurrección).
  \item La asignación especie--entorno es defendible y la $\mu_{\mathrm{opt}}$
  relativa tiene sentido ($2{,}1 > 0{,}26 > 0{,}069$~h$^{-1}$).
\end{itemize}

\end{document}
